\documentclass[12pt, a4paper]{article}
\usepackage[T2A]{fontenc}
\usepackage[utf8]{inputenc}
\usepackage[russian]{babel}
\usepackage{amsmath, amssymb}
\usepackage{enumitem}
\usepackage{geometry}
\geometry{top=2cm, bottom=2cm, left=2.5cm, right=2cm}

\begin{document}

\section*{Георгий Игоревич Вольфсон <<Делимость с человеческим лицом>>}

\subsection*{Беседа четвёртая. Десятичная запись числа.}

\subsubsection*{Упражнения}

\begin{enumerate}[label=\arabic*., wide=0pt, leftmargin=*]

\item \textbf{Найдите все трехзначные числа, состоящие из различных цифр, все три слова в названии которых начинаются с одной и той же буквы (например, для двухзначного числа подошло бы "двадцать девять").}

\textbf{Решение:} 
Составим таблицу первых букв:
\begin{itemize}
\item Единицы: о (один), д (два), т (три), ч (четыре), п (пять), ш (шесть), с (семь), в (восемь), д (девять).
\item Десятки (для чисел от 20 до 99, так как 11–19 дают только два слова в названии трёхзначного числа): д (двадцать), т (тридцать), с (сорок), п (пятьдесят), ш (шестьдесят), с (семьдесят), в (восемьдесят), д (девяносто).
\item Сотни: с (сто), д (двести), т (триста), ч (четыреста), п (пятьсот), ш (шестьсот), с (семьсот), в (восемьсот), д (девятьсот).
\end{itemize}
Единственная буква, которая встречается во всех трёх наборах, — это <<с>> (семь, сорок/семьдесят, сто/семьсот). Подходит вариант: \textbf{сто сорок семь} — число 147. 
Проверка: цифры различны, все три слова начинаются на <<с>>. Других вариантов нет. 
Ответ: 147.

\item \textbf{Найдите все натуральные числа, большие своей последней цифры ровно в 6 раз.}

\textbf{Решение:}
Пусть число имеет вид \(10a + b\), где \(b\) — последняя цифра. Условие:
\[
10a + b = 6b \quad\Rightarrow\quad 10a = 5b \quad\Rightarrow\quad 2a = b.
\]
\(a\) — цифра от 1 до 9, \(b\) — цифра от 0 до 9, причём \(b\) чётная. Перебираем:
\begin{itemize}
\item \(b = 2 \Rightarrow a = 1\) \(\Rightarrow\) число 12,
\item \(b = 4 \Rightarrow a = 2\) \(\Rightarrow\) число 24,
\item \(b = 6 \Rightarrow a = 3\) \(\Rightarrow\) число 36,
\item \(b = 8 \Rightarrow a = 4\) \(\Rightarrow\) число 48.
\end{itemize}
При \(b = 0\) получаем \(a = 0\), что не является натуральным числом.
Ответ: 12, 24, 36, 48.

\item \textbf{Из двухзначного числа вычли число, записанное теми же цифрами, но в обратном порядке. Докажите, что полученная разность кратна 9.}

\textbf{Решение:}
Запишем исходное число как \(\overline{ab} = 10a + b\), а число с обратным порядком цифр как \(\overline{ba} = 10b + a\). Тогда
\[
\overline{ab} - \overline{ba} = (10a + b) - (10b + a) = 9a - 9b = 9(a - b).
\]
Полученная разность содержит множитель 9, следовательно, она всегда делится на 9.

\item \textbf{К трехзначному числу приписали справа его же. Докажите, что полученное число делится на 13.}

\textbf{Решение:}
Пусть исходное трёхзначное число равно \(\overline{abc} = 100a + 10b + c\). Приписав его справа, получим шестизначное число:
\[
\overline{abcabc} = 1000 \cdot \overline{abc} + \overline{abc} = 1001 \cdot \overline{abc}.
\]
Заметим, что \(1001 = 7 \cdot 11 \cdot 13\). Следовательно,
\[
\overline{abcabc} = 13 \cdot (77 \cdot \overline{abc}).
\]
Таким образом, полученное число всегда кратно 13.

\item \textbf{Найдите двухзначное число, обладающее следующим свойством: если зачеркнуть его последнюю цифру, то получится число, в 14 раз меньшее.}

\textbf{Решение:}
Пусть число имеет вид \(\overline{ab} = 10a + b\). Зачёркиваем последнюю цифру, остаётся \(a\). Условие:
\[
10a + b = 14a \quad\Rightarrow\quad b = 4a.
\]
Поскольку \(a\) и \(b\) — цифры (\(1 \le a \le 9\), \(0 \le b \le 9\)), возможны варианты:
\begin{itemize}
\item \(a = 1 \Rightarrow b = 4\) \(\Rightarrow\) число 14. \(14 = 14 \cdot 1\) — верно.
\item \(a = 2 \Rightarrow b = 8\) \(\Rightarrow\) число 28. \(28 = 14 \cdot 2\) — верно.
\end{itemize}
Ответ: 14 и 28.

\item \textbf{Шестизначное число начинается с цифры 2. Откинув эту цифру слева и приписав её справа, получим число, которое в 3 раза больше первоначального. Найдите первоначальное число.}

\textbf{Решение:}
Пусть исходное число имеет вид \(2 \cdot 10^5 + n\), где \(n\) — пятизначное число (от 00000 до 99999). После перестановки цифры 2 в конец получим число \(10n + 2\). Условие:
\[
10n + 2 = 3 \cdot (2 \cdot 10^5 + n).
\]
Решаем:
\[
10n + 2 = 600000 + 3n,
\]
\[
7n = 599998,
\]
\[
n = 85714.
\]
Исходное число: \(285714\). 
Проверка: откидываем 2 слева — 85714, приписываем 2 справа — 857142. Действительно, \(857142 = 3 \cdot 285714\).

\item \textbf{Сумма цифр искомого двухзначного числа равна 8. Если цифры переставить, то получится число, которое меньше на 18. Чему равно искомое число?}

\textbf{Решение:}
Пусть число равно \(\overline{ab} = 10a + b\). Условия:
\[
a + b = 8,
\]
\[
10a + b = 10b + a + 18.
\]
Из второго уравнения:
\[
9a - 9b = 18 \quad\Rightarrow\quad a - b = 2.
\]
Решаем систему:
\[
\begin{cases}
a + b = 8,\\
a - b = 2,
\end{cases}
\Rightarrow
2a = 10 \Rightarrow a = 5,\; b = 3.
\]
Искомое число: 53. 
Проверка: 53, перестановка — 35, разность \(53 - 35 = 18\).

\item \textbf{Найдите двухзначное число, зная, что число его единиц на 4 меньше числа его десятков, а произведение искомого числа и суммы его цифр равно 1330.}

\textbf{Решение:}
Пусть число равно \(\overline{ab} = 10a + b\). Условия:
\[
a = b + 4,
\]
\[
(10a + b)(a + b) = 1330.
\]
Подставляем \(a = b + 4\):
\[
(10(b+4) + b)((b+4) + b) = (10b + 40 + b)(2b + 4) = (11b + 40)(2b + 4) = 1330.
\]
Раскрываем скобки:
\[
(11b + 40) \cdot 2(b + 2) = 1330 \quad\Rightarrow\quad 2(11b + 40)(b + 2) = 1330,
\]
\[
(11b + 40)(b + 2) = 665.
\]
\[
11b^2 + 22b + 40b + 80 = 665 \quad\Rightarrow\quad 11b^2 + 62b + 80 - 665 = 0,
\]
\[
11b^2 + 62b - 585 = 0.
\]
Дискриминант:
\[
D = 62^2 - 4 \cdot 11 \cdot (-585) = 3844 + 25740 = 29584 = 172^2.
\]
\[
b = \frac{-62 + 172}{2 \cdot 11} = \frac{110}{22} = 5 \quad (\text{отрицательный корень не подходит}).
\]
Тогда \(a = b + 4 = 9\). Искомое число: 95. 
Проверка: \(95\), сумма цифр \(14\), произведение \(95 \cdot 14 = 1330\).

\item \textbf{После деления некоторого двухзначного числа на сумму его цифр в частном получается 7 и в остатке 6. После деления этого же числа на произведение его цифр в частном получается 3 и в остатке 11. Найдите это двухзначное число.}

\textbf{Решение:}
Пусть число равно \(\overline{ab} = 10a + b\). Условия:
\[
10a + b = 7(a + b) + 6,
\]
\[
10a + b = 3ab + 11.
\]
Из первого уравнения:
\[
10a + b = 7a + 7b + 6 \quad\Rightarrow\quad 3a - 6b = 6 \quad\Rightarrow\quad a - 2b = 2 \quad\Rightarrow\quad a = 2b + 2.
\]
Подставляем во второе:
\[
10(2b + 2) + b = 3b(2b + 2) + 11,
\]
\[
20b + 20 + b = 6b^2 + 6b + 11,
\]
\[
21b + 20 = 6b^2 + 6b + 11,
\]
\[
0 = 6b^2 + 6b + 11 - 21b - 20 = 6b^2 - 15b - 9.
\]
Делим на 3:
\[
2b^2 - 5b - 3 = 0.
\]
Дискриминант:
\[
D = (-5)^2 - 4 \cdot 2 \cdot (-3) = 25 + 24 = 49,\quad \sqrt{D} = 7.
\]
\[
b = \frac{5 + 7}{2 \cdot 2} = \frac{12}{4} = 3 \quad (\text{отрицательный корень не подходит}).
\]
Тогда \(a = 2 \cdot 3 + 2 = 8\). Искомое число: 83. 
Проверка: сумма цифр \(11\), \(83 = 7 \cdot 11 + 6\) (верно). Произведение цифр \(24\), \(83 = 3 \cdot 24 + 11\) (верно).

\item \textbf{Найдите число, которое оканчивается на 2, а при перестановке этой цифры в начало уменьшается в два раза.}

\textbf{Решение:}
Пусть исходное число имеет вид \(10x + 2\), где \(x\) — некоторое натуральное число или ноль (часть числа без последней цифры). После перестановки цифры 2 в начало получим число \(2 \cdot 10^k + x\), где \(k\) — количество цифр в числе \(x\) (то есть \(10^k \le x < 10^{k+1}\)). Условие:
\[
10x + 2 = 2 \cdot (2 \cdot 10^k + x).
\]
Раскрываем скобки:
\[
10x + 2 = 4 \cdot 10^k + 2x,
\]
\[
8x = 4 \cdot 10^k - 2,
\]
\[
4x = 2 \cdot 10^k - 1.
\]
Левая часть \(4x\) — чётное число. Правая часть \(2 \cdot 10^k - 1\) — нечётное (так как \(2 \cdot 10^k\) чётно, вычитаем 1, получаем нечётное). Чётное число не может равняться нечётному. Следовательно, решений нет. 
Ответ: таких чисел не существует.

\end{enumerate}

\end{document}